%% ========================================================================
%% Chapter 2: Manajemen Perangkat Keras dan Perintah Dasar
%% ========================================================================

\chapter{Manajemen Perangkat Keras dan Perintah Dasar}
\label{ch:hardware-and-basic-commands}

Bayangkan kamu baru membeli laptop baru dan menyalakannya untuk pertama kali. Sistem operasi perlu tahu apa saja komponen di dalamnya: jenis prosesor, kapasitas RAM, kartu grafis, dan disk. Di Windows, informasi ini tersedia di Device Manager. Di Linux, kamu mendapatkan akses yang jauh lebih dalam lewat terminal.

Linux mendeteksi perangkat keras melalui proses bertahap. Saat komputer dinyalakan, kernel melakukan pemindaian terhadap semua bus seperti PCI dan USB, lalu mencocokkan perangkat yang ditemukan dengan modul kernel yang sesuai. Hasil deteksi ini bisa ditelusuri melalui perintah seperti \cmd{lscpu}, \cmd{lspci}, \cmd{lsusb}, dan \cmd{dmesg}. Prosesnya berlangsung otomatis, tetapi jejaknya tetap bisa kamu audit dari terminal.

Memahami perangkat keras dan perintah dasar terminal adalah fondasi dari semua praktikum Linux berikutnya. Tanpa keterampilan ini, kamu tidak akan bisa membaca informasi sistem, berpindah direktori dengan percaya diri, mengelola berkas sederhana, atau menafsirkan dokumentasi perintah yang kamu temui sepanjang buku ini.

\textbf{Yang Akan Kamu Pelajari}

Setelah menyelesaikan bab ini, kamu akan mampu:

\begin{enumerate}
    \item mengidentifikasi komponen perangkat keras menggunakan perintah \cmd{lspci}, \cmd{lsusb}, \cmd{lscpu}, dan \cmd{dmesg};
    \item menjelaskan fungsi modul kernel dan melihat modul yang aktif dengan \cmd{lsmod} dan \cmd{modinfo};
    \item menggunakan perintah dasar terminal untuk navigasi serta operasi file dan direktori;
    \item membaca isi berkas dan output perintah dengan \cmd{cat}, \cmd{head}, dan \cmd{tail};
    \item menyaring dan mengolah teks menggunakan \cmd{grep}, \cmd{sed}, dan \cmd{awk};
    \item menggunakan dokumentasi bawaan seperti \cmd{man}, \cmd{--help}, dan \cmd{info} untuk mempelajari perintah baru secara mandiri.
\end{enumerate}

\begin{notebox}
Seluruh praktikum pada bab ini dijalankan pada direktori \dir{~/lab-os/week02-hardware/}.
Buat direktori ini sebelum memulai praktikum pertama:
\begin{lstlisting}[language=bash]
mkdir -p ~/lab-os/week02-hardware
cd ~/lab-os/week02-hardware
\end{lstlisting}
\end{notebox}

%% ------------------------------------------------------------------------
\section{Deteksi Perangkat Keras}
\label{sec:hardware-detection}

Setiap kali Linux menyala, kernel memeriksa semua perangkat yang terpasang. Hasilnya tersimpan dan bisa kamu baca kapan saja. Perintah \cmd{lscpu} menampilkan arsitektur prosesor, jumlah inti fisik, jumlah thread, dan fitur virtualisasi. Perintah \cmd{lspci} memperlihatkan semua perangkat di bus PCI, termasuk kartu jaringan, kartu grafis, dan pengendali penyimpanan. Untuk perangkat USB, gunakan \cmd{lsusb}. Jika kamu ingin laporan perangkat keras yang lebih lengkap dalam satu perintah, \cmd{lshw} adalah pilihan yang tepat meski memerlukan hak akses superuser.

Pesan kernel yang dihasilkan saat boot tersimpan dalam buffer cincin yang bisa dibaca dengan \cmd{dmesg}. Di sana kamu bisa melihat catatan deteksi perangkat: kapan disk ditemukan, driver apa yang dimuat, dan apakah ada perangkat yang gagal diinisialisasi. Tambahkan flag \texttt{-T} untuk menampilkan waktu dalam format yang bisa dibaca manusia.

\begin{tipbox}
Beberapa perintah deteksi perangkat keras mungkin belum terpasang di instalasi minimal. Pasang paket yang diperlukan dengan:
\begin{lstlisting}[language=bash]
sudo apt update
sudo apt install -y pciutils usbutils lshw
\end{lstlisting}
\end{tipbox}

\subsection*{Praktikum 2.1: Identifikasi Komponen Perangkat Keras VM}

Tujuan praktikum ini adalah membaca spesifikasi prosesor, memori, dan perangkat PCI dari mesin virtual.

\begin{enumerate}
    \item Tampilkan informasi prosesor:
\begin{lstlisting}[language=bash, caption={Membaca spesifikasi prosesor}]
lscpu
\end{lstlisting}
    \textit{Amati:} Perhatikan baris \texttt{CPU(s)}, \texttt{Core(s) per socket}, dan \texttt{Vendor ID}. Mesin virtual biasanya memperlihatkan prosesor host dengan jumlah inti yang dibatasi.

    \item Lihat daftar perangkat PCI:
\begin{lstlisting}[language=bash, caption={Daftar perangkat PCI beserta driver aktif}]
lspci -nnk
\end{lstlisting}
    \textit{Amati:} Setiap entri menampilkan ID vendor dan perangkat dalam tanda kurung, diikuti baris \texttt{Kernel driver in use} yang menunjukkan nama modul driver aktif.

    \item Cari kartu jaringan dan driver yang menanganinya:
\begin{lstlisting}[language=bash, caption={Filter informasi kartu jaringan}]
lspci -nnk | grep -A3 -i ethernet
\end{lstlisting}

    \item Lihat pesan kernel dari proses boot:
\begin{lstlisting}[language=bash, caption={Pesan kernel saat boot}]
dmesg -T | head -50
\end{lstlisting}
    \textit{Amati:} Cari baris yang menyebutkan nama perangkat seperti \texttt{sda}, \texttt{eth0}, atau \texttt{ens3}. Baris itu menandai saat kernel pertama kali mendeteksi perangkat tersebut.

    \item Cari pesan boot yang berkaitan dengan virtualisasi atau perangkat penyimpanan:
\begin{lstlisting}[language=bash, caption={Menyaring pesan kernel dengan dmesg}]
dmesg -T | grep -Ei "virtio|nvme|sata|usb"
\end{lstlisting}
\end{enumerate}

\begin{challengebox}
Pilih satu modul kernel yang aktif dari output \cmd{lsmod}. Gunakan \cmd{modinfo} untuk mencari deskripsi, lisensi, dan nama berkas modulnya. Lalu gunakan \cmd{dmesg -T | grep -i nama-modul} untuk memeriksa apakah ada pesan kernel yang berkaitan dengan modul tersebut saat sistem menyala.
\end{challengebox}

%% ------------------------------------------------------------------------
\section{Modul Kernel dan Driver Perangkat}
\label{sec:kernel-modules}

Kernel Linux tidak memuat semua kode driver sekaligus. Sebaliknya, driver tersedia sebagai \textit{modul kernel}, yaitu berkas kode yang bisa dimuat saat dibutuhkan dan dilepas ketika tidak lagi diperlukan. Prinsip ini sama dengan menancapkan steker alat listrik hanya saat akan digunakan, bukan membiarkannya terus terpasang. Modul tersimpan di \dir{/lib/modules/\$(uname -r)/} dan biasanya memiliki ekstensi \texttt{.ko}.

Perintah \cmd{lsmod} memperlihatkan semua modul yang sedang aktif bersama ukuran dan jumlah proses yang menggunakannya. Untuk melihat detail satu modul, gunakan \cmd{modinfo nama-modul}: kamu akan melihat deskripsi, parameter yang tersedia, nama penulis, lisensi, dan path file modulnya. Perintah \cmd{modprobe nama-modul} memuat modul beserta semua dependensinya secara otomatis.

Saat kernel berhasil memuat modul, hasilnya biasanya terlihat di pesan \cmd{dmesg}, di daftar perangkat PCI/USB, atau pada kemampuan baru yang tersedia di sistem. Karena itu, tiga alat yang paling sering dipakai bersama adalah \cmd{lsmod} untuk melihat apa yang aktif, \cmd{modinfo} untuk memahami modul tertentu, dan \cmd{dmesg} untuk menelusuri kapan driver tersebut dipakai atau menimbulkan error.

\begin{warningbox}
Jangan melepas modul yang sedang digunakan oleh sistem aktif. Melepas driver disk atau modul jaringan dapat menyebabkan sistem tidak responsif atau koneksi SSH terputus.
\end{warningbox}

\subsection*{Praktikum 2.2: Eksplorasi Modul Aktif dan Driver Perangkat}

Tujuan praktikum ini adalah melihat modul kernel yang sedang berjalan dan menghubungkannya dengan perangkat yang terdeteksi sistem.

\begin{enumerate}
    \item Tampilkan versi kernel yang sedang berjalan:
\begin{lstlisting}[language=bash, caption={Versi kernel aktif}]
uname -r
\end{lstlisting}

    \item Lihat daftar modul aktif:
\begin{lstlisting}[language=bash, caption={Daftar modul kernel aktif}]
lsmod | head -20
\end{lstlisting}
    \textit{Amati:} Kolom pertama adalah nama modul, kolom kedua ukuran dalam byte, kolom ketiga jumlah proses yang menggunakan modul ini.

    \item Baca detail modul \texttt{loop} (modul untuk membuat disk virtual dari berkas):
\begin{lstlisting}[language=bash, caption={Detail modul loop}]
modinfo loop
\end{lstlisting}

    \item Pastikan modul \texttt{loop} termuat, lalu verifikasi:
\begin{lstlisting}[language=bash, caption={Muat dan verifikasi modul loop}]
sudo modprobe loop
lsmod | grep loop
\end{lstlisting}

    \item Lihat bagaimana \cmd{modprobe} menyelesaikan dependensi modul tanpa benar-benar memuat ulang modul:
\begin{lstlisting}[language=bash, caption={Simulasi pemuatan modul dan dependensinya}]
modprobe -n -v loop
\end{lstlisting}
    \textit{Amati:} Flag \texttt{-n} melakukan simulasi, sedangkan \texttt{-v} menampilkan langkah yang akan diambil. Ini berguna ketika kamu ingin memahami dependensi modul tanpa mengubah keadaan sistem.

    \item Hubungkan modul dengan perangkat nyata menggunakan keluaran PCI dan pesan kernel:
\begin{lstlisting}[language=bash, caption={Hubungkan modul, perangkat, dan pesan kernel}]
lspci -nnk | grep -A3 -i -E "ethernet|network|vga|storage"
dmesg -T | grep -Ei "virtio|nvme|usb|ethernet" | tail -20
\end{lstlisting}
\end{enumerate}

\begin{challengebox}
Pilih satu perangkat dari output \cmd{lspci -nnk}. Catat nama perangkat dan driver aktifnya. Lalu gunakan \cmd{modinfo} untuk mencari deskripsi modul tersebut, kemudian verifikasi dengan \cmd{dmesg} apakah ada pesan boot yang menyebut perangkat atau driver yang sama.
\end{challengebox}

%% ------------------------------------------------------------------------
\section{Perintah Dasar Terminal}
\label{sec:basic-commands}

Terminal Linux adalah antarmuka utama administrasi server. Tidak seperti antarmuka grafis yang menyembunyikan detail, terminal memberi kontrol penuh: kamu tahu persis apa yang terjadi dan mengapa. Penguasaan perintah dasar bukan pilihan, melainkan keharusan.

Navigasi dimulai dengan tiga perintah: \cmd{pwd} menampilkan lokasi direktori saat ini, \cmd{ls} memperlihatkan isi direktori, dan \cmd{cd} berpindah ke direktori lain. Flag \texttt{-l} pada \cmd{ls} menampilkan format panjang dengan permission, pemilik, ukuran, dan waktu modifikasi. Flag \texttt{-a} memperlihatkan berkas tersembunyi yang diawali titik. Flag \texttt{-h} menyajikan ukuran dalam satuan yang mudah dibaca seperti KB atau MB.

Membuat dan mengelola berkas menggunakan \cmd{mkdir} untuk membuat direktori, \cmd{touch} untuk membuat berkas kosong, \cmd{cp} untuk menyalin, \cmd{mv} untuk memindahkan atau mengganti nama, dan \cmd{rm} untuk menghapus. Flag \texttt{-p} pada \cmd{mkdir} membuat seluruh jalur direktori sekaligus, termasuk direktori induk yang belum ada. Flag \texttt{-r} pada \cmd{cp} dan \cmd{rm} menerapkan operasi secara rekursif ke isi direktori.

Untuk membaca isi berkas, \cmd{cat} menampilkan semua konten sekaligus, \cmd{head} menampilkan beberapa baris pertama, dan \cmd{tail} menampilkan beberapa baris terakhir. Flag \texttt{-n} pada keduanya menentukan jumlah baris.

\begin{warningbox}
Perintah \cmd{rm -rf} menghapus direktori beserta seluruh isinya tanpa konfirmasi. Tidak ada tempat sampah. Verifikasi path yang akan dihapus sebelum menekan Enter.
\end{warningbox}

\subsection*{Praktikum 2.3: Buat Struktur Direktori Latihan}

Tujuan praktikum ini adalah membangun struktur direktori kerja dan melatih manipulasi berkas dasar.

\begin{enumerate}
    \item Pindah ke direktori kerja bab ini:
\begin{lstlisting}[language=bash, caption={Masuk ke direktori kerja}]
cd ~/lab-os/week02-hardware
pwd
\end{lstlisting}

    \item Buat struktur direktori bertingkat sekaligus:
\begin{lstlisting}[language=bash, caption={Membuat struktur direktori}]
mkdir -p proyek/src proyek/docs proyek/logs
ls -l proyek/
\end{lstlisting}
    \textit{Amati:} Flag \texttt{-p} membuat semua direktori dalam jalur tersebut. Tanpa \texttt{-p}, perintah akan gagal jika direktori induk belum ada.

    \item Buat beberapa berkas contoh dan isi salah satunya:
\begin{lstlisting}[language=bash, caption={Membuat dan mengisi berkas}]
touch proyek/docs/readme.txt
echo "INFO: layanan dimulai" > proyek/logs/app.log
echo "WARN: disk hampir penuh" >> proyek/logs/app.log
echo "ERROR: koneksi gagal" >> proyek/logs/app.log
cat proyek/logs/app.log
\end{lstlisting}

    \item Salin berkas log ke direktori lain:
\begin{lstlisting}[language=bash, caption={Menyalin berkas}]
cp proyek/logs/app.log proyek/docs/app-backup.log
ls -lh proyek/docs/
\end{lstlisting}

    \item Ganti nama berkas backup:
\begin{lstlisting}[language=bash, caption={Mengganti nama berkas}]
mv proyek/docs/app-backup.log proyek/docs/app-2024.log
ls proyek/docs/
\end{lstlisting}

    \item Tampilkan tiga baris terakhir dari log:
\begin{lstlisting}[language=bash, caption={Membaca akhir berkas log}]
tail -3 proyek/logs/app.log
\end{lstlisting}
\end{enumerate}

\begin{challengebox}
Buat struktur direktori \texttt{arsip/2024/01}, \texttt{arsip/2024/02}, dan \texttt{arsip/2024/03} menggunakan \textbf{satu perintah} \cmd{mkdir -p} dengan ekspansi brace. Periksa hasilnya dengan \cmd{ls -R arsip/}. Kemudian hapus seluruh direktori \texttt{arsip} sekaligus menggunakan \cmd{rm -r arsip/}.
\end{challengebox}

%% ------------------------------------------------------------------------
\section{Manipulasi Teks}
\label{sec:text-manipulation}

Hampir semua data di Linux berbentuk teks: log sistem, berkas konfigurasi, output perintah. Kemampuan menyaring, mengubah, dan mengekstrak informasi dari teks adalah keterampilan yang akan kamu gunakan setiap hari sebagai administrator sistem.

Perintah \cmd{grep} mencari baris yang cocok dengan pola tertentu. Flag \texttt{-n} menampilkan nomor baris, \texttt{-i} mengabaikan perbedaan huruf besar-kecil, \texttt{-v} menampilkan baris yang tidak cocok, dan \texttt{-r} melakukan pencarian rekursif di semua berkas dalam direktori. Untuk mencari dua pola sekaligus, gunakan flag \texttt{-E} dengan pola \texttt{pola1|pola2}.

Perintah \cmd{sed} adalah editor stream yang memproses teks baris per baris. Operasi yang paling sering digunakan adalah substitusi dengan sintaks \texttt{s/lama/baru/g}: huruf \texttt{s} berarti substitusi, \texttt{g} di akhir berarti ganti semua kemunculan dalam satu baris. Flag \texttt{-i} menerapkan perubahan langsung ke berkas asli. Untuk keamanan saat bekerja di sistem produksi, gunakan \texttt{-i.bak} agar berkas asli disimpan sebagai cadangan.

Perintah \cmd{awk} memproses teks berdasarkan kolom. Setiap baris dibagi menjadi kolom yang dipisahkan spasi, dan kamu bisa mengakses kolom tertentu dengan \texttt{\$1}, \texttt{\$2}, dan seterusnya. Variabel \texttt{\$NF} merujuk ke kolom terakhir. Variabel \texttt{NR} berisi nomor baris saat ini. Kamu bisa menambahkan kondisi seperti \texttt{NR>1} untuk melewati baris header.

\subsection*{Praktikum 2.4: Filter dan Ekstrak Teks dari Output Sistem}

Tujuan praktikum ini adalah menggunakan \cmd{grep}, \cmd{sed}, dan \cmd{awk} untuk mengolah output perintah sistem nyata.

\begin{enumerate}
    \item Siapkan berkas log latihan dengan lebih banyak konten:
\begin{lstlisting}[language=bash, caption={Membuat berkas log latihan}]
cat > ~/lab-os/week02-hardware/sistem.log << 'EOF'
INFO: layanan web dimulai pada port 8080
WARN: penggunaan CPU mencapai 75%
ERROR: gagal membaca berkas konfigurasi
INFO: cadangan data selesai
WARN: memori tersisa kurang dari 20%
ERROR: koneksi database terputus
INFO: pembaruan sistem dijadwalkan
EOF
cat sistem.log
\end{lstlisting}

    \item Cari baris yang berisi kata ERROR:
\begin{lstlisting}[language=bash, caption={Filter baris ERROR}]
grep "ERROR" sistem.log
\end{lstlisting}
    \textit{Amati:} Hasilnya hanya baris dengan kata ERROR. Coba tambahkan flag \texttt{-n} untuk melihat nomor baris.

    \item Tampilkan baris yang bukan INFO:
\begin{lstlisting}[language=bash, caption={Tampilkan baris selain INFO}]
grep -v "INFO" sistem.log
\end{lstlisting}

    \item Ganti kata \texttt{WARN} menjadi \texttt{PERINGATAN} di semua baris:
\begin{lstlisting}[language=bash, caption={Substitusi teks dengan sed}]
sed 's/WARN/PERINGATAN/g' sistem.log
\end{lstlisting}
    \textit{Amati:} Berkas asli tidak berubah. Hanya output yang tampil di terminal yang sudah dimodifikasi.

    \item Buat berkas konfigurasi latihan dan ubah nilainya:
\begin{lstlisting}[language=bash, caption={Edit konfigurasi dengan sed}]
cat > server.conf << 'EOF'
PORT=8080
MODE=development
HOST=localhost
EOF
sed -i 's/development/production/g' server.conf
cat server.conf
\end{lstlisting}

    \item Ekstrak kolom pertama dan terakhir dari output \cmd{df -h}:
\begin{lstlisting}[language=bash, caption={Ekstrak kolom dengan awk}]
df -h | awk '{print $1, $NF}'
\end{lstlisting}
    \textit{Amati:} \texttt{\$1} adalah kolom pertama (nama filesystem), \texttt{\$NF} adalah kolom terakhir (titik kait). Output jauh lebih ringkas dari output \cmd{df -h} penuh.

    \item Filter filesystem yang penggunaan disk di atas 50\%:
\begin{lstlisting}[language=bash, caption={Filter dengan kondisi awk}]
df -h | awk 'NR==1 {print} NR>1 && ($5+0) > 50 {print}'
\end{lstlisting}
\end{enumerate}

\begin{challengebox}
Gabungkan \cmd{grep} dan \cmd{awk} dalam satu pipeline untuk mengekstrak informasi dari \file{/proc/meminfo}. Cari baris yang mengandung kata \texttt{Mem} menggunakan \cmd{grep}, lalu ekstrak nama field (kolom pertama) dan nilainya (kolom kedua) menggunakan \cmd{awk}. Perintahnya: \texttt{grep Mem /proc/meminfo | awk '\{print \$1, \$2, \$3\}'}. Tambahkan \cmd{sort} di akhir pipeline untuk mengurutkan hasilnya.
\end{challengebox}

%% ------------------------------------------------------------------------
\section{Bantuan dan Dokumentasi Perintah}
\label{sec:command-documentation}

Jumlah perintah di Linux sangat banyak, dan tidak ada administrator yang menghafal semuanya. Keterampilan yang lebih penting adalah mengetahui cara mencari jawaban dengan cepat. Linux menyediakan dokumentasi bawaan yang sangat kaya: halaman manual, ringkasan opsi lewat \cmd{--help}, dokumentasi GNU lewat \cmd{info}, dan indeks topik melalui \cmd{apropos}.

Perintah \cmd{man nama-perintah} menampilkan halaman manual lengkap berisi deskripsi, opsi, contoh, dan berkas terkait. Hampir semua utilitas modern juga mendukung \cmd{nama-perintah --help} untuk ringkasan cepat. Untuk beberapa program GNU, \cmd{info nama-perintah} menyediakan dokumentasi yang lebih panjang dan terstruktur. Sementara itu, \cmd{apropos kata-kunci} berguna ketika kamu belum tahu nama perintahnya, misalnya mencari semua perintah yang berkaitan dengan ``copy'' atau ``archive''.

Ada juga perintah bantu lain yang sering dipakai di terminal: \cmd{which} untuk menemukan lokasi executable di \var{PATH}, \cmd{type} untuk mengetahui apakah sebuah nama adalah builtin shell, alias, fungsi, atau executable biasa, serta \cmd{help} untuk bantuan perintah bawaan shell Bash.

Pendekatan ini penting untuk buku ini: alih-alih berhenti ketika menjumpai perintah baru, biasakan membuka dokumentasinya, membaca opsi yang relevan, lalu mencoba contoh kecil secara langsung di terminal.

\subsection*{Praktikum 2.5: Telusuri Dokumentasi Perintah}

Tujuan praktikum ini adalah membiasakan diri mencari dokumentasi dan memahami perbedaan sumber bantuan yang tersedia di Linux.

\begin{enumerate}
    \item Buka halaman manual untuk \cmd{ls} dan cari bagian yang menjelaskan flag \texttt{-a}:
\begin{lstlisting}[language=bash, caption={Membuka halaman manual}]
man ls
\end{lstlisting}
    \textit{Amati:} Di dalam \cmd{man}, gunakan tombol \texttt{/} untuk mencari kata kunci seperti \texttt{-a}, lalu tekan \texttt{n} untuk hasil berikutnya dan \texttt{q} untuk keluar.

    \item Lihat ringkasan opsi cepat untuk \cmd{grep} dan \cmd{tar}:
\begin{lstlisting}[language=bash, caption={Ringkasan bantuan cepat}]
grep --help | head -20
tar --help | head -20
\end{lstlisting}

    \item Gunakan \cmd{apropos} untuk mencari perintah berdasarkan topik:
\begin{lstlisting}[language=bash, caption={Mencari perintah berdasarkan kata kunci}]
apropos copy
apropos archive
\end{lstlisting}

    \item Periksa lokasi dan tipe perintah \cmd{cd}, \cmd{ls}, dan \cmd{grep}:
\begin{lstlisting}[language=bash, caption={Mengenali jenis perintah di shell}]
type cd
type ls
which grep
\end{lstlisting}

    \item Baca bantuan untuk builtin shell dan dokumentasi info bila tersedia:
\begin{lstlisting}[language=bash, caption={Bantuan builtin dan info}]
help cd
info coreutils 'ls invocation'
\end{lstlisting}
\end{enumerate}

\begin{challengebox}
Pilih tiga perintah dari bab ini, misalnya \cmd{ls}, \cmd{grep}, dan \cmd{tar}. Untuk masing-masing, catat: satu cara membuka dokumentasi lengkap, satu cara melihat bantuan singkat, dan satu opsi yang menurutmu paling berguna. Jelaskan kapan kamu akan memakai opsi tersebut.
\end{challengebox}

%% ------------------------------------------------------------------------
\section{Rangkuman}

Dalam bab ini, kamu telah mempelajari:

\begin{itemize}[noitemsep,topsep=2pt]
    \item \textbf{Deteksi perangkat keras}: perintah \cmd{lspci}, \cmd{lsusb}, \cmd{lscpu}, dan \cmd{dmesg} memperlihatkan komponen sistem yang dideteksi oleh kernel.
    \item \textbf{Modul kernel}: \cmd{lsmod} menampilkan modul aktif, \cmd{modinfo} membaca detail modul, dan \cmd{modprobe} memuat modul beserta dependensinya.
    \item \textbf{Perintah dasar terminal}: \cmd{ls}, \cmd{cd}, \cmd{pwd}, \cmd{mkdir -p}, \cmd{cp}, \cmd{mv}, \cmd{rm}, \cmd{cat}, \cmd{head}, dan \cmd{tail} adalah fondasi navigasi dan manipulasi berkas.
    \item \textbf{Manipulasi teks}: \cmd{grep} menyaring baris berdasarkan pola, \cmd{sed} mengganti teks, \cmd{awk} mengekstrak kolom dari output berformat.
    \item \textbf{Dokumentasi perintah}: \cmd{man}, \cmd{--help}, \cmd{info}, \cmd{apropos}, \cmd{type}, dan \cmd{which} membantu kamu mempelajari perintah baru secara mandiri.
\end{itemize}

%% ------------------------------------------------------------------------
\section{Tugas Praktikum}

\textbf{Instruksi Umum}: Kerjakan seluruh tugas secara mandiri. Dokumentasikan setiap langkah dengan tangkapan layar dan kumpulkan dalam format PDF.

\begin{exercisebox}[2.1 Profil Perangkat Keras VM]
Jalankan \cmd{lspci -nnk} dan pilih satu perangkat PCI. Catat: (a) nama perangkat, (b) ID vendor dan perangkat dalam format heksadesimal, (c) nama driver kernel aktif. Kemudian jalankan \cmd{modinfo} untuk driver tersebut dan tuliskan deskripsi singkat serta nama penulis modulnya.
\end{exercisebox}

\begin{exercisebox}[2.2 Dokumentasi dan Modul Kernel]
Pilih satu modul aktif dari hasil \cmd{lsmod}. Dokumentasikan: (a) nama modul, (b) deskripsi dan lisensi dari \cmd{modinfo}, (c) satu pesan kernel terkait dari \cmd{dmesg} bila ada. Setelah itu, gunakan \cmd{man} atau \cmd{--help} untuk mencari arti salah satu opsi dari perintah \cmd{modprobe} dan jelaskan fungsinya.
\end{exercisebox}

\begin{exercisebox}[2.3 Analisis Log dengan grep/sed/awk]
Buat berkas \file{server.log} berisi minimal 10 baris dengan variasi tingkat pesan: INFO, WARN, dan ERROR. Kemudian: (a) gunakan \cmd{grep -c ERROR server.log} untuk menghitung jumlah baris ERROR, (b) gunakan \cmd{sed} untuk mengganti semua kata \texttt{WARN} menjadi \texttt{WARNING} dan simpan ke berkas baru \file{server-clean.log}, (c) gunakan \cmd{awk} untuk menampilkan hanya kata pertama (tingkat pesan) dari setiap baris. Sertakan tangkapan layar output setiap perintah.
\end{exercisebox}

%% ------------------------------------------------------------------------
\section{Referensi}

\begin{itemize}
    \item Shotts, W. \textit{The Linux Command Line}, edisi ke-2. No Starch Press, 2019.
    \item Nemeth, E., Snyder, G., Hein, T. R. \textit{UNIX and Linux System Administration Handbook}, edisi ke-5. Pearson, 2017.
    \item Ward, B. \textit{How Linux Works}, edisi ke-3. No Starch Press, 2021.
    \item Dokumentasi Ubuntu Server: \url{https://ubuntu.com/server/docs}
    \item Kernel.org Documentation: \url{https://www.kernel.org/doc/html/latest/}
\end{itemize}
